\section{Introduzione}
\subsection{Applicazioni client-server}
Le applicazioni Internet si basano su un’architettura client-server, cioè sono composte da due parti software logicamente distinte:
\begin{itemize}
    \item \textbf{Lato server}: fornisce il servizio
    \item \textbf{Lato client}: richiede il servizio
\end{itemize}
Un server può gestire più client contemporaneamente. L'utente, invece, accede al servizio attraverso un user agent, ovvero un software 
che unisce il client con l’interfaccia grafica usata dalla persona.

\begin{figure}[ht]
    \centering
    \includegraphics[width=0.6\textwidth]{img/client_server.png}
\end{figure}

Il funzionamento delle applicazioni client-server è il seguente:
\begin{itemize}
    \item Il server è collegato ad internet e attende le richieste dei client.
    \item Il client si collega ad internet per richiedere un servizio al server.
    \item La comunicazione avviene grazie a un protocollo applicativo, cioè un insieme di regole di scambio dei messaggi.
    \item L’utente interagisce tramite un user agent, che rappresenta l’interfaccia grafica e la parte client dell’applicazione.
\end{itemize}

Ogni applicazione client-server si basa su un protocollo, cioè un insieme di regole che stabiliscono come client e server devono comunicare, 
va progettato con precisione fin dall’inizio, perché rappresenta la base di tutto il funzionamento dell’applicazione.
\newline
È importante chiarire alcuni concetti fondamentali:
\begin{itemize}
    \item Client e server non sono macchine, ma programmi. Lo stesso computer può eseguire sia un client che un server.
    \item I server di solito stanno su computer più potenti, perché devono gestire tante richieste contemporaneamente.
    \item I client invece possono essere leggeri e funzionare su molti tipi di dispositivi, dai PC agli smartphone.
    \item Nulla vieta che un computer faccia da server e allo stesso tempo ospiti anche uno o più client della stessa applicazione.
\end{itemize}

\subsection{Layer Stack}
Internet funziona a strati (stack di protocolli). Le applicazioni (come posta, web, ecc.) non comunicano direttamente con la rete fisica, ma si 
appoggiano al livello di trasporto (es. TCP o UDP).

\begin{figure}[ht]
    \centering
    \includegraphics[width=0.6\textwidth]{img/pila_protocollare.png}
\end{figure}

\subsubsection{Transport Layer}
Per permettere la comunicazione tra programmi, il livello di trasporto usa le porte:
\begin{itemize}
    \item Ogni programma (client o server) che comunica in rete deve “aprire” una porta numerata.
    \item L’indirizzo IP identifica il computer, mentre la porta identifica il programma che gira su quel computer.
    \item I server hanno porte fisse, su cui restano in ascolto (per esempio, il web server usa la porta 80 per HTTP).
    \item I client di solito usano porte temporanee, assegnate automaticamente dal sistema.
    \item Le porte da 0 a 1023 sono già riservate ad applicazioni ben note (HTTP, FTP, SMTP, ecc.).
\end{itemize}

Le applicazioni in internet possono scegliere tra due protocolli di trasporto:
\begin{itemize}
    \item \textbf{TCP (Transmission Control Protocol)}: protocollo con connessione avanzata, è come una telefonata.
            Prima ci si accorda (handshake), poi si parla in modo ordinato e sicuro. Assicura che i dati arrivino corretti e nell’ordine giusto.
    \item \textbf{UDP (User Datagram Protocol)}: protocollo senza connessione, è come spedire una lettera senza raccomandata.
            Non c’è garanzia che arrivi, né che arrivi in ordine. È più veloce e leggero, usato per applicazioni in tempo reale (es.\ videochiamate, giochi online).
\end{itemize}

Vediamo ora un esempio di comunicazione TCP.\
\begin{figure}[ht]
    \centering
    \includegraphics[width=0.8\textwidth]{img/TCP_1.png}
\end{figure}
\newline
Il server rimane in ascolto su una porta, ad esempio la 80 per HTTP, pronto a ricevere richieste dai client.

\begin{figure}[ht]
    \centering
    \includegraphics[width=0.8\textwidth]{img/TCP_2.png}
\end{figure}
Il client invia una richiesta di connessione al server, specificando l’indirizzo IP del server e la porta su cui è in ascolto. Il server accetta la connessione
e stabilisce un canale di comunicazione dedicato con quel client.
\begin{figure}[ht]
    \centering
    \includegraphics[width=0.8\textwidth]{img/TCP_3.png}
\end{figure}

Da questo momento, server e client comunicano attraverso questo canale, mentre il server rimane libero di ascoltare nuove richieste da altri client.
\begin{figure}[ht]
    \centering
    \includegraphics[width=0.8\textwidth]{img/TCP_4.png}
\end{figure}

\newpage
\subsection{Socket}
Un \textbf{socket} è come una porta virtuale che permette al client e al server di collegarsi al livello di trasporto.
\begin{itemize}
    \item Ogni socket è associato a una porta.
    \item Client e server creano un canale di comunicazione virtuale attraverso internet.
\end{itemize}

\begin{figure}[ht]
    \centering
    \includegraphics[width=0.6\textwidth]{img/socket.png}
\end{figure}

Un socket è implementato come una struttura dati e memorizza, tra le altre, queste informazioni:
\begin{itemize}
    \item \textbf{Tipo di socket}: \begin{itemize}
        \item \textbf{Stream socket}: usato con TCP, garantisce una comunicazione affidabile e ordinata.
        \item \textbf{Datagram socket}: usato con UDP, è più veloce ma non garantisce l’ordine né l’affidabilità.
    \end{itemize}
    \item \textbf{Indirizzo locale}: IP dell'host + porta del processo locale.
    \item \textbf{Indirizzo remoto}: IP dell'altro host + porta del processo remoto.
\end{itemize}

\newpage
Vediamo ora come funzionano i socket e in che modo client e server comunicano tra loro.

\begin{figure}[ht]
    \centering
    \includegraphics[width=0.8\textwidth]{img/s_1.png}
\end{figure}
\begin{figure}[ht]
    \centering
    \includegraphics[width=0.8\textwidth]{img/s_2.png}
\end{figure}
\begin{figure}[ht]
    \centering
    \includegraphics[width=0.8\textwidth]{img/s_3.png}
\end{figure}

\newpage

Con TCP, il server ha una porta specifica (porta applicativa) su cui rimane in ascolto. Quando un client vuole connettersi, invia una richiesta a quella porta.
Quando arriva la richiesta di un client, il server crea un nuovo socket attivo dedicato a quella connessione. In questo modo, la comunicazione con quel client 
non blocca l’ascolto di nuove richieste.\newline
Un server può avere quindi:
\begin{itemize}
    \item un socket passivo, per ascoltare le richieste in arrivo.
    \item più socket attivi, uno per ogni connessione stabilita con un client.
\end{itemize}


Il server ascolta su un socket passivo, collegato a una porta applicativa nota. Per ogni nuovo client, il server apre un socket attivo dedicato, ovviamente se due 
client girano sulla stessa macchina, useranno porte diverse.
\newline
Una connessione TCP è identificata in modo univoco da 4 elementi:
\texttt{<IP server, porta server, IP client, porta client>}

\subsection{Elementi chiave nella progettazione client-server}
Per progettare un’applicazione client-server su Internet bisogna fare alcune scelte fondamentali:
\begin{itemize}
    \item \textbf{Definire il servizio}: spiegare bene che cosa fa l’applicazione (es.\ inviare email, trasferire file, ecc.).
    \item \textbf{Definire il protocollo}: stabilire con precisione le regole di comunicazione tra client e server.
    \item \textbf{Scegliere il trasporto}: decidere se usare TCP (affidabile, ordinato) o UDP (veloce, non garantito).
    \item \textbf{Scegliere la porta}: scegliere su quale porta il server resterà in ascolto per accettare le connessioni.
\end{itemize}

Quando un server usa TCP, può essere progettato in due modi diversi:
\begin{itemize}
    \item \textbf{Modello sincrono}: ogni richiesta viene servita in modo sequenziale o con un thread dedicato -$>$ più facile da programmare, ma può richiedere molte risorse.
    \item \textbf{Modello asincrono}: il server usa un unico thread per gestire più connessioni contemporaneamente -$>$ più complesso, ma più efficiente in termini di risorse.
\end{itemize}
Per capire meglio la differenza, si usa spesso la metafora del ristorante, dove i clienti rappresentano i client e i camerieri rappresentano il server.\newline
Per spiegare i due modelli di server TCP (sincrono e asincrono) possiamo immaginare un ristorante:
\begin{flushleft}
    \textbf{Il servizio}
    \begin{itemize}
        \item Quando un cliente arriva, il cameriere lo accoglie, prende l’ordine, lo porta in cucina, e poi serve il piatto.
        \item Alcune operazioni (come aspettare che la cucina prepari il piatto) possono bloccare il cameriere.
    \end{itemize}
\end{flushleft}

\begin{enumerate}
    \item \textbf{Modello sincrono}\begin{itemize}
        \item Ogni cliente ha un cameriere dedicato.
        \item Se un cameriere aspetta l’ordine dalla cucina, gli altri continuano a servire i loro clienti.
        \item È semplice da gestire, ma servono tante risorse (molti camerieri = molti thread).
    \end{itemize}
    \item \textbf{Modello asincrono}\begin{itemize}
        \item Tutti i clienti vengono serviti da un unico cameriere.
        \item Se il cameriere rimanesse fermo in attesa della cucina, bloccherebbe tutti gli altri clienti.
        \item Per evitare questo, il cameriere lavora in modo non bloccante: prende gli ordini, li lascia in cucina, passa ad altri clienti e torna solo quando i piatti sono pronti.
        \item Il cameriere usa una coda di eventi per gestire più clienti contemporaneamente.
    \end{itemize}
\end{enumerate}

\subsubsection{Sincrono vs Asincrono}
\begin{itemize}
    \item Il modello sincrono è più semplice da programmare, ma se i client aumentano richiede tanti thread → quindi può diventare poco scalabile.
    \item Il modello asincrono è più difficile da programmare, ma è altamente scalabile → spesso basta un solo thread per gestire tutti i client.
\end{itemize}

Alcuni sistemi combinano i due approcci:
\begin{itemize}
    \item \textbf{Thread pool} → un numero limitato di thread che gestiscono i client.
    \item \textbf{Gestione asincrona delle operazioni bloccanti} → ogni thread può delegare i compiti che rischiano di bloccarlo.
\end{itemize}
Una novità è rappresentata dai \textbf{virtual threads} di Java, che offrono la semplicità della programmazione sincrona con la scalabilità tipica delle architetture asincrone.